% Supplementary Note 4: penalty-strength sweep, design and coverage; replication across training seeds, design.
% Sources: run lists and records of the sweep analysis (D/sweep_analysis_final_20261105, commit 8595a16; paper/figure_data), docs/PREREG.md (history),
% fit ledgers under D/tier12_v3 and D/jhpce_tier12, experiments/sweep_replication/ (branch sweep-replication, a475ea3). Numbers: paper/supp-analysis_numbers.csv.
\par\medskip\noindent\textbf{Supplementary Note 4. Penalty-strength sweep: design and coverage.}\par\smallskip
\begingroup\small\setlength{\parskip}{3pt}
\noindent\textit{Design.} Plain scVI and the six penalties (Supplementary Note 1) were fitted on all eight tasks at the ten half-decade strengths 0.1, 0.3, 1, 3, 10, 30, 100, 300, 1,000 and 3,000, with the batch-conditioned and the unconditioned decoder and with seeds 100, 101 and 102. Edge strengths were added in two places: the discriminator with the unconditioned decoder was also fitted at 0.01 and 0.03, and the Sinkhorn divergence with the batch-conditioned decoder at 10,000 and 30,000. Each task therefore has 378 fits: 6 of plain scVI (three seeds, two decoders), 360 of the penalties (six penalties, ten strengths, three seeds, two decoders) and 12 at edge strengths (two places, two strengths, three seeds), so the sweep has 3,024 fits on the eight tasks, 96 of them at edge strengths. Nine further fits, at strength 0.03 for the discriminator, the reference critic and the pooled critic on the immune task with the batch-conditioned decoder (seeds 100--102), come from an earlier strength-extension round (below). The analysis reads every fitted strength, so it reads 3,033 fits, 387 of them on the immune task. A penalty's curve therefore has ten strengths, twelve for the discriminator with the unconditioned decoder and for the Sinkhorn divergence with the batch-conditioned decoder on every task, and 11 for the discriminator, the reference critic and the pooled critic on the immune task with the batch-conditioned decoder. None of the 3,033 fits failed (Supplementary Table S5).

\noindent\textit{Why the edge strengths exist.} The sweep was first run as a pilot on the three tasks fitted first (the small chromatin-accessibility, immune and simulation 1 tasks). A strength-choice rule, written before the pilot ran, was to read the pilot's seed means and set, for each penalty family and decoder, a range of six strengths for a planned fixed-strength benchmark; where a range reached the end of the grid, the rule asked for two more half-decade strengths beyond that end. On the three tasks this happened in two places, the discriminator with the unconditioned decoder at the low end and the Sinkhorn divergence with the batch-conditioned decoder at the high end, and the two extra strengths were fitted there (36 fits on the three tasks). When the sweep was extended to the other five tasks, the same 378 fits per task were made, the edge strengths included, so every task has the same strengths. The nine further fits come from a second rule, which matched batch removal at a strength nearest a target batch score and, where that strength was the lowest of its grid, fitted the next lower half-decade value; this happened for three penalties on the immune task with the batch-conditioned decoder. They are not strengths of the strength-choice rule, but they are fitted points of these penalties' curves and are read like every other strength. That pair is low-signal and is not ranked (Supplementary Note 5).

\noindent\textit{The planned fixed-strength benchmark was not run.} The strength-choice rules were written to set up a benchmark at six strengths per penalty family and decoder with five new seeds. That benchmark was not run, so the strength-choice rules choose no strength: every analysis reads every sweep strength, and the comparison is exploratory. The replication across training seeds (below) refits strengths chosen from the sweep's batch scores, not from its biology scores.

\noindent\textit{Machines.} All fits of a task were made on one kind of machine, so that the hardware never differs between penalties, strengths or decoders within a task: the small chromatin-accessibility, immune and simulation 1 tasks on the workstation's NVIDIA RTX 3080 graphics card, and each of the other five tasks in one cluster job on NVIDIA L40S graphics cards (Supplementary Note 11).

\noindent\textit{Sinkhorn solver.} All 528 Sinkhorn fits of the sweep (66 per task) were made with the corrected Sinkhorn solver, at commit 10a1ca5 (Supplementary Note 1 gives the solver). On the three workstation tasks, the fits made with the earlier solver were replaced: the 180 fits at the ten grid strengths were refitted on 17--18 October 2026 and the 18 fits at the edge strengths 10,000 and 30,000 on 19 October 2026, 198 in all. The five cluster tasks were fitted from 22 October 2026 at the same commit.

\noindent\textit{Replication across training seeds: design.} For each of the 13 ranked task and decoder pairs and each of the six penalties, we took the two sweep strengths that bracket the first crossing of $b^*$ on the sweep's seed-mean curve, with $b^*$ fixed at its sweep value (Supplementary Note 5), and refitted them with five new seeds (0--4). 46 of the 78 curves cross $b^*$ more than once; the first crossing sets the bracket, as in the sweep reading. In one cell, the discriminator with the batch-conditioned decoder on the human--mouse immune task, the curve already lies above $b^*$ at the lowest sweep strength, 0.1, so the bracket runs from plain scVI to that strength and plain scVI at the new seeds is its lower point. Plain scVI was refitted at seeds 0--4 on all eight tasks with both decoders. This gives 855 fits, 775 with a penalty and 80 of plain scVI. The three low-signal pairs (the immune task and both simulations with the batch-conditioned decoder) were skipped. If the new seed means at the two strengths do not bracket $b^*$, the next sweep strength outward is fitted once with seeds 0--4; if $b^*$ is still not bracketed, the penalty takes the bottom rank in that pair and is marked as not replicating the crossing. The fits of this extension step are counted in Supplementary Note 11, and the readings that needed it are marked in Supplementary Table S9.

\noindent\textit{Replication: provenance.} The design script \texttt{experiments/sweep\_replication/replication\_design.py} wrote the run list (\texttt{manifests/sweep\_replication\_x1.tsv}, 855 rows) and its summary in \texttt{experiments/sweep\_replication/design\_20261102/} (commit 2113ec7). The prediction file \texttt{PREDICTION.json} (the sweep readings and ranks at $b^*$ of the 13 ranked pairs and the 4 within-pair differences of 0.04 or more) and the agreement rule \texttt{AGREEMENT.md} were committed under the tag \texttt{replication-prediction-v1} (commit 3a90d0f) before any replication fit was scored. The prediction file records the table it was read from, the ranking table of the sweep analysis at commit a361ba2, with its SHA-256 (786cb8c8\ldots); that table is byte-identical to the ranking table of the final sweep analysis used in this paper (\texttt{paper/figure\_data/ranking\_pairs.csv}), so the prediction holds the readings and ranks of Supplementary Table S6.
\par\endgroup
